\section{Technical choices}
\label{sec:task3}

\subsection{Low / ultra-low power sensor choice}
    First we wanted to find a sensor that can fit with our RF Energy harvesting. 
The challenge is that the power given by our circuit is very low, we could expect around 3V
and few µA. By making some researches on electronic components resellers website, we can find 
some interesting low / ultra-low power sensors. In a first time, we selected three components, a voltage supervisor (TPS3839), 
a temperature sensor (STTS22H) and a capacitive touch and sensitivity switch (used for DAS, PCF8883).
By analysing their datasheets   

La puissance captée par notre antenne réceptrice dépend directement de la distance et des gains des antennes, et peut être modélisée par l'équation des télécommunications de Friis (Équation \ref{eq:friis}) :

\begin{equation} \label{eq:friis}
    P_{r} = P_{t} + G_{t} + G_{r} + 20\log_{10}\left(\frac{\lambda}{4\pi d}\right)
\end{equation}

Où $P_r$ est la puissance reçue en dBm, $\lambda$ la longueur d'onde en mètres, et $d$ la distance de la source en mètres.

\subsection{Efficacité de conversion RF-DC}
L'élément critique du système est la \textit{rectenna} (composée de l'antenne et du circuit redresseur). Le rendement global $\eta$ dépend fortement de la puissance RF incidente à l'entrée de la diode Schottky :

\begin{equation} \label{eq:rendement}
    \eta = \frac{P_{DC}}{P_{RF\_in}} \times 100
\end{equation} efficency 